\documentclass[11pt]{article}
\usepackage[margin=1in]{geometry}
\usepackage{xcolor}
\usepackage[breakable,listings]{tcolorbox}
\tcbuselibrary{listings}
\usepackage{hyperref}
\usepackage{enumitem}
\usepackage{listings}
\usepackage{verbatim}
\usepackage{amsmath}
\usepackage{graphicx}
\usepackage{booktabs}
\usepackage{float}

% Define colors
\definecolor{osaorange}{RGB}{220,115,38}
\definecolor{aiblue}{RGB}{30,144,255}
\definecolor{codebackground}{RGB}{248,248,248}

% Configure hyperref
\hypersetup{
    colorlinks=true,
    linkcolor=blue,
    urlcolor=blue,
    citecolor=blue
}

% Configure listings for code display
\lstset{
  basicstyle=\small\ttfamily,
  breaklines=true,
  breakatwhitespace=true,
  postbreak=\mbox{\textcolor{gray}{$\hookrightarrow$}\space},
  columns=flexible,
  keepspaces=true,
  showstringspaces=false
}

% Code environment using tcolorbox with listings
\newtcblisting{rcode}{
  colback=white,
  colframe=gray!50,
  boxrule=0.5pt,
  arc=0pt,
  left=3pt,
  right=3pt,
  top=3pt,
  bottom=3pt,
  width=\textwidth,
  breakable,
  listing only,
  listing options={
    basicstyle=\small\ttfamily,
    breaklines=true,
    breakatwhitespace=true,
    postbreak=\mbox{\textcolor{gray}{$\hookrightarrow$}\space},
    columns=flexible,
    keepspaces=true,
    showstringspaces=false
  }
}

\title{}
\author{}
\date{}

\begin{document}

% Header box
\begin{tcolorbox}[colback=osaorange,colframe=black,boxrule=2pt,arc=0pt,left=3pt,right=3pt,top=6pt,bottom=6pt,boxsep=2pt]
\centering
\textcolor{white}{\textbf{\Large H524 Introduction to Biostatistics}}\\
\textcolor{white}{\textbf{\Large Week 9 Lab: Correlation and Simple Linear Regression}}\\
\textcolor{white}{\textbf{\Large Fall 2025}}
\end{tcolorbox}

\vspace{8pt}

\noindent\textbf{Format:} Online\\
\textbf{Tools Required:} R, RStudio\\
\textbf{Estimated Time:} 90 minutes

\section{Learning Objectives}

By the end of this lab, you will be able to:
\begin{enumerate}
\item Calculate and interpret Pearson and Spearman correlation coefficients
\item Test the significance of correlation coefficients
\item Create scatter plots with informative labels and regression lines
\item Fit simple linear regression models using \texttt{lm()}
\item Interpret regression output (coefficients, $R^2$, $p$-values)
\item Check regression assumptions using diagnostic plots
\item Make predictions using regression models
\item Calculate confidence and prediction intervals
\item \textcolor{aiblue}{\textbf{Use AI tools to help with regression analysis and interpretation}}
\end{enumerate}

\section{Part 1: Correlation Analysis}

Correlation measures the strength and direction of the linear relationship between two continuous variables.

\subsection{Exercise 1.1: Calculating Correlation}

\textbf{Scenario:} We have data on age and systolic blood pressure for 15 adults.

\begin{rcode}
# Age and systolic blood pressure data
age <- c(35, 38, 40, 42, 45, 47, 48, 50, 52, 55, 58, 60, 62, 65, 68)
sbp <- c(118, 122, 125, 128, 130, 132, 135, 136, 138, 140, 145, 148,
         150, 155, 158)

# Calculate Pearson correlation coefficient
cor_result <- cor.test(age, sbp)
print(cor_result)

# Extract key information
cat("\nPearson correlation coefficient: r =", round(cor_result$estimate, 3), "\n")
cat("95% Confidence Interval:", round(cor_result$conf.int, 3), "\n")
cat("t-statistic:", round(cor_result$statistic, 3), "\n")
cat("p-value:", format.pval(cor_result$p.value, digits=3), "\n")

# Interpretation
if(cor_result$p.value < 0.05) {
  cat("\nConclusion: There is significant evidence of correlation (p < 0.05)\n")
} else {
  cat("\nConclusion: No significant evidence of correlation (p >= 0.05)\n")
}
\end{rcode}

\textbf{Questions:}
\begin{enumerate}[label=\alph*)]
\item What is the value of the correlation coefficient? Is it positive or negative?
\item What does this tell us about the relationship between age and blood pressure?
\item Is the correlation statistically significant?
\end{enumerate}

\vspace{1.5cm}

\subsection{Exercise 1.2: Visualizing Correlation}

Always visualize your data before computing statistics!

\begin{rcode}
# Create scatter plot
plot(age, sbp,
     main = "Systolic Blood Pressure vs. Age",
     xlab = "Age (years)",
     ylab = "Systolic BP (mmHg)",
     pch = 19,
     col = "blue",
     cex = 1.2)

# Add grid for easier reading
grid()

# Add correlation coefficient to plot
text(40, 155,
     paste("r =", round(cor(age, sbp), 3)),
     pos = 4, cex = 1.2, col = "red")

# Optionally, add a lowess smoother to see the trend
lines(lowess(age, sbp), col = "red", lwd = 2)
\end{rcode}

\textbf{Observation:} Does the scatter plot support the correlation coefficient you calculated? Do you see a linear trend?

\vspace{1cm}

\subsection{Exercise 1.3: Spearman Correlation}

Spearman correlation is based on ranks and is more robust to outliers.

\begin{rcode}
# Calculate Spearman correlation
spearman_result <- cor.test(age, sbp, method = "spearman")
print(spearman_result)

# Compare Pearson and Spearman
pearson_r <- cor(age, sbp, method = "pearson")
spearman_r <- cor(age, sbp, method = "spearman")

cat("\nPearson r:", round(pearson_r, 3), "\n")
cat("Spearman rho:", round(spearman_r, 3), "\n")
cat("Difference:", round(abs(pearson_r - spearman_r), 3), "\n")
\end{rcode}

\textbf{Question:} Why are Pearson and Spearman correlations similar in this case? When would they differ substantially?

\vspace{1.5cm}

\subsection{Exercise 1.4: Effect of Outliers on Correlation}

Let's see how outliers affect correlation.

\begin{rcode}
# Create data with an outlier
age_outlier <- c(age, 45)
sbp_outlier <- c(sbp, 95)  # Unusually low BP for age 45

# Calculate correlations with and without outlier
r_original <- cor(age, sbp)
r_outlier <- cor(age_outlier, sbp_outlier)

cat("Original correlation:", round(r_original, 3), "\n")
cat("With outlier:", round(r_outlier, 3), "\n")
cat("Change:", round(r_outlier - r_original, 3), "\n")

# Visualize the effect
par(mfrow = c(1, 2))

# Original data
plot(age, sbp, main = "Original Data",
     xlab = "Age", ylab = "SBP", pch = 19, col = "blue")
text(40, 155, paste("r =", round(r_original, 3)), col = "red")

# With outlier
plot(age_outlier, sbp_outlier, main = "With Outlier",
     xlab = "Age", ylab = "SBP", pch = 19, col = "blue")
points(45, 95, col = "red", pch = 19, cex = 2)  # Highlight outlier
text(40, 155, paste("r =", round(r_outlier, 3)), col = "red")

par(mfrow = c(1, 1))
\end{rcode}

\textbf{Key lesson:} Always examine your data visually. A single outlier can dramatically change the correlation!

\vspace{0.5cm}

\section{Part 2: Simple Linear Regression}

Regression allows us to model the relationship and make predictions.

\subsection{Exercise 2.1: Fitting a Regression Model}

\begin{rcode}
# Fit simple linear regression: sbp ~ age
model <- lm(sbp ~ age)

# Display full summary
summary(model)

# Extract key components
coef_table <- summary(model)$coefficients
r_squared <- summary(model)$r.squared
adj_r_squared <- summary(model)$adj.r.squared
residual_se <- summary(model)$sigma

cat("\n=== Regression Results Summary ===\n")
cat("Intercept (b0):", round(coef(model)[1], 2), "mmHg\n")
cat("Slope (b1):", round(coef(model)[2], 2), "mmHg/year\n")
cat("R-squared:", round(r_squared, 4), "\n")
cat("Adjusted R-squared:", round(adj_r_squared, 4), "\n")
cat("Residual standard error:", round(residual_se, 2), "mmHg\n")

# Regression equation
cat("\nRegression equation:\n")
cat("SBP =", round(coef(model)[1], 2), "+",
    round(coef(model)[2], 2), "* Age\n")
\end{rcode}

\textbf{Interpretation questions:}
\begin{enumerate}[label=\alph*)]
\item What is the slope? Interpret it in context.
\item What is the intercept? Is it meaningful in this context?
\item What percentage of variability in SBP is explained by age?
\item Is the relationship statistically significant?
\end{enumerate}

\vspace{2cm}

\subsection{Exercise 2.2: Visualizing the Regression Line}

\begin{rcode}
# Scatter plot with regression line
plot(age, sbp,
     main = "Systolic BP vs. Age with Regression Line",
     xlab = "Age (years)",
     ylab = "Systolic BP (mmHg)",
     pch = 19,
     col = "blue",
     cex = 1.2)

# Add regression line
abline(model, col = "red", lwd = 2)

# Add confidence band
age_pred <- seq(min(age), max(age), length.out = 100)
conf_interval <- predict(model,
                         newdata = data.frame(age = age_pred),
                         interval = "confidence",
                         level = 0.95)

# Plot confidence bands
lines(age_pred, conf_interval[,"lwr"], col = "gray", lty = 2, lwd = 1.5)
lines(age_pred, conf_interval[,"upr"], col = "gray", lty = 2, lwd = 1.5)

# Add equation to plot
equation <- paste("SBP =", round(coef(model)[1], 1), "+",
                  round(coef(model)[2], 2), "* Age")
text(40, 155, equation, pos = 4, col = "red", cex = 1)

# Add R-squared
text(40, 150, paste("R-squared =", round(r_squared, 3)),
     pos = 4, col = "red", cex = 1)

# Add legend
legend("topleft",
       legend = c("Observed data", "Regression line", "95% CI"),
       col = c("blue", "red", "gray"),
       lty = c(NA, 1, 2),
       pch = c(19, NA, NA),
       lwd = c(NA, 2, 1.5))
\end{rcode}

\subsection{Exercise 2.3: Making Predictions}

\begin{rcode}
# Predict SBP for specific ages
new_ages <- data.frame(age = c(45, 55, 65))

# Point predictions
predictions <- predict(model, newdata = new_ages)
cat("Point predictions:\n")
for(i in 1:nrow(new_ages)) {
  cat("Age", new_ages$age[i], "years: SBP =",
      round(predictions[i], 1), "mmHg\n")
}

# Confidence intervals for mean response
conf_int <- predict(model, newdata = new_ages,
                    interval = "confidence", level = 0.95)
cat("\n95% Confidence Intervals (for mean SBP):\n")
print(round(conf_int, 1))

# Prediction intervals for individual response
pred_int <- predict(model, newdata = new_ages,
                    interval = "prediction", level = 0.95)
cat("\n95% Prediction Intervals (for individual SBP):\n")
print(round(pred_int, 1))

# Compare widths
cat("\nNotice: Prediction intervals are WIDER than confidence intervals\n")
cat("CI width for age 55:", round(conf_int[2,"upr"] - conf_int[2,"lwr"], 1), "mmHg\n")
cat("PI width for age 55:", round(pred_int[2,"upr"] - pred_int[2,"lwr"], 1), "mmHg\n")
\end{rcode}

\textbf{Question:} Why are prediction intervals wider than confidence intervals?

\vspace{1.5cm}

\subsection{Exercise 2.4: Checking Regression Assumptions}

Regression requires several assumptions. Let's check them!

\begin{rcode}
# Create 2x2 panel of diagnostic plots
par(mfrow = c(2, 2))
plot(model)
par(mfrow = c(1, 1))

# Residuals vs Fitted: checks linearity and constant variance
# Normal Q-Q: checks normality of residuals
# Scale-Location: checks homoscedasticity
# Residuals vs Leverage: checks for influential points

# Additional tests
cat("\n=== Assumption Checks ===\n")

# Test for normality of residuals
shapiro_test <- shapiro.test(residuals(model))
cat("\nShapiro-Wilk test for normality:\n")
cat("W =", round(shapiro_test$statistic, 4),
    ", p-value =", round(shapiro_test$p.value, 4), "\n")
if(shapiro_test$p.value > 0.05) {
  cat("Residuals appear normally distributed (p > 0.05)\n")
} else {
  cat("Evidence against normality (p < 0.05)\n")
}

# Examine residuals
cat("\nResidual summary:\n")
print(summary(residuals(model)))
\end{rcode}

\textbf{What to look for:}
\begin{itemize}
\item \textbf{Residuals vs Fitted:} Should show random scatter around zero (no pattern)
\item \textbf{Normal Q-Q:} Points should follow diagonal line (normality)
\item \textbf{Scale-Location:} Should show random scatter (constant variance)
\item \textbf{Residuals vs Leverage:} No points with high leverage and large residuals
\end{itemize}

\vspace{0.5cm}

\subsection{Exercise 2.5: ANOVA Table for Regression}

\begin{rcode}
# ANOVA table
anova_table <- anova(model)
print(anova_table)

# Extract components
ss_total <- sum((sbp - mean(sbp))^2)
ss_regression <- anova_table$`Sum Sq`[1]
ss_residual <- anova_table$`Sum Sq`[2]

cat("\n=== ANOVA Decomposition ===\n")
cat("Total SS:", round(ss_total, 2), "\n")
cat("Regression SS:", round(ss_regression, 2), "\n")
cat("Residual SS:", round(ss_residual, 2), "\n")
cat("R-squared = Regression SS / Total SS =",
    round(ss_regression / ss_total, 4), "\n")

# F-test
f_stat <- anova_table$`F value`[1]
p_value <- anova_table$`Pr(>F)`[1]
cat("\nF-statistic:", round(f_stat, 2), "\n")
cat("p-value:", format.pval(p_value, digits=3), "\n")
cat("\nConclusion: The regression model is",
    ifelse(p_value < 0.05, "significant", "not significant"), "\n")
\end{rcode}

\section{Part 3: Real-World Examples}

\subsection{Exercise 3.1: Birth Weight and Gestational Age}

\begin{rcode}
# Simulated data based on realistic values
set.seed(524)
gestational_age <- round(runif(50, 32, 42))
birth_weight <- 150 * gestational_age - 2800 + rnorm(50, 0, 250)

# Create data frame
birth_data <- data.frame(weeks = gestational_age, weight = birth_weight)

# Summary
cat("Birth Weight Study\n")
cat("Sample size:", nrow(birth_data), "\n")
cat("Gestational age range:", range(birth_data$weeks), "weeks\n")
cat("Birth weight range:", round(range(birth_data$weight)), "grams\n\n")

# Correlation
cor_birth <- cor.test(birth_data$weeks, birth_data$weight)
cat("Pearson correlation: r =", round(cor_birth$estimate, 3), "\n")
cat("p-value:", format.pval(cor_birth$p.value, digits=3), "\n\n")

# Regression
model_birth <- lm(weight ~ weeks, data = birth_data)
summary(model_birth)

# Plot
plot(birth_data$weeks, birth_data$weight,
     main = "Birth Weight vs. Gestational Age",
     xlab = "Gestational Age (weeks)",
     ylab = "Birth Weight (grams)",
     pch = 19, col = "darkgreen", cex = 1.2)
abline(model_birth, col = "red", lwd = 2)

# Add interpretation
cat("\n=== Clinical Interpretation ===\n")
cat("Slope:", round(coef(model_birth)[2], 1), "grams per week\n")
cat("Each additional week of gestation increases birth weight by",
    round(coef(model_birth)[2], 0), "grams on average\n")
cat("R-squared:", round(summary(model_birth)$r.squared, 3), "\n")
cat("Gestational age explains",
    round(100 * summary(model_birth)$r.squared, 1),
    "% of variability in birth weight\n")
\end{rcode}

\textbf{Practice:} Use the model to predict birth weight for babies at 36, 38, and 40 weeks gestation. Include 95\% prediction intervals.

\vspace{2cm}

\subsection{Exercise 3.2: Exercise and Resting Heart Rate}

\begin{rcode}
# Simulated data
set.seed(123)
exercise_sessions <- sample(0:7, 30, replace = TRUE)
resting_hr <- 78 - 2.5 * exercise_sessions + rnorm(30, 0, 5)

# Create data frame
hr_data <- data.frame(sessions = exercise_sessions, hr = resting_hr)

# Scatter plot
plot(hr_data$sessions, hr_data$hr,
     main = "Resting Heart Rate vs. Exercise Frequency",
     xlab = "Exercise Sessions per Week",
     ylab = "Resting Heart Rate (bpm)",
     pch = 19, col = "purple", cex = 1.2)

# Correlation
cor_hr <- cor.test(hr_data$sessions, hr_data$hr)
cat("Correlation: r =", round(cor_hr$estimate, 3), "\n")
cat("p-value:", format.pval(cor_hr$p.value, digits=3), "\n\n")

# Regression
model_hr <- lm(hr ~ sessions, data = hr_data)
summary(model_hr)
abline(model_hr, col = "red", lwd = 2)

# Interpretation
cat("\n=== Interpretation ===\n")
cat("Slope:", round(coef(model_hr)[2], 2), "bpm per session\n")
cat("Each additional exercise session per week is associated with",
    abs(round(coef(model_hr)[2], 1)), "bpm decrease in resting heart rate\n")

# Check assumptions
par(mfrow = c(2, 2))
plot(model_hr)
par(mfrow = c(1, 1))
\end{rcode}

\section{Part 4: Advanced Topics}

\subsection{Exercise 4.1: Comparing Two Regression Models}

\begin{rcode}
# Compare models with different predictors
# Example: Height and weight data
set.seed(456)
height_cm <- rnorm(40, 170, 10)
weight_kg <- 0.8 * height_cm - 70 + rnorm(40, 0, 5)
age_years <- sample(20:60, 40, replace = TRUE)

# Model 1: Weight predicted by height
model1 <- lm(weight_kg ~ height_cm)

# Model 2: Weight predicted by age
model2 <- lm(weight_kg ~ age_years)

# Compare R-squared
cat("Model 1 (Height): R-squared =", round(summary(model1)$r.squared, 3), "\n")
cat("Model 2 (Age): R-squared =", round(summary(model2)$r.squared, 3), "\n\n")

# Compare AIC (lower is better)
cat("Model 1 (Height): AIC =", round(AIC(model1), 2), "\n")
cat("Model 2 (Age): AIC =", round(AIC(model2), 2), "\n\n")

# Visualize both models
par(mfrow = c(1, 2))
plot(height_cm, weight_kg, main = "Model 1: Height",
     xlab = "Height (cm)", ylab = "Weight (kg)", pch = 19)
abline(model1, col = "red", lwd = 2)
text(160, 100, paste("R² =", round(summary(model1)$r.squared, 3)))

plot(age_years, weight_kg, main = "Model 2: Age",
     xlab = "Age (years)", ylab = "Weight (kg)", pch = 19)
abline(model2, col = "blue", lwd = 2)
text(25, 100, paste("R² =", round(summary(model2)$r.squared, 3)))
par(mfrow = c(1, 1))
\end{rcode}

\subsection{Exercise 4.2: Influence of Sample Size}

\begin{rcode}
# Demonstrate effect of sample size on inference
set.seed(789)

sample_sizes <- c(10, 30, 100, 300)
results <- data.frame(n = sample_sizes, r = NA, p_value = NA,
                      se_slope = NA, ci_width = NA)

for(i in 1:length(sample_sizes)) {
  n <- sample_sizes[i]

  # Generate data
  x <- rnorm(n, 50, 10)
  y <- 2 + 0.5 * x + rnorm(n, 0, 5)

  # Correlation
  cor_result <- cor.test(x, y)

  # Regression
  model_temp <- lm(y ~ x)
  se_slope <- summary(model_temp)$coefficients[2, "Std. Error"]
  ci <- confint(model_temp)[2, ]
  ci_width <- ci[2] - ci[1]

  # Store results
  results$r[i] <- cor_result$estimate
  results$p_value[i] <- cor_result$p.value
  results$se_slope[i] <- se_slope
  results$ci_width[i] <- ci_width
}

cat("Effect of Sample Size on Inference:\n\n")
print(round(results, 4))

cat("\nKey observation: As sample size increases,\n")
cat("  - Standard errors decrease\n")
cat("  - Confidence intervals narrow\n")
cat("  - Power to detect relationships increases\n")
\end{rcode}

\section{Practice Problems}

\subsection{Problem 1: Cholesterol and Age}

A study examines the relationship between age and total cholesterol in 25 adults.

\begin{rcode}
# Data
set.seed(111)
age_chol <- seq(30, 70, length.out = 25)
cholesterol <- 160 + 1.2 * age_chol + rnorm(25, 0, 15)
\end{rcode}

\textbf{Tasks:}
\begin{enumerate}[label=\alph*)]
\item Create a scatter plot of cholesterol vs. age
\item Calculate the Pearson correlation and test its significance
\item Fit a linear regression model
\item Interpret the slope in context
\item Check regression assumptions
\item Predict cholesterol level for a 50-year-old with 95\% prediction interval
\end{enumerate}

\vspace{2cm}

\subsection{Problem 2: BMI and Blood Pressure}

Examine the relationship between BMI and systolic blood pressure.

\begin{rcode}
# Data
set.seed(222)
bmi <- rnorm(35, 27, 5)
sbp_bmi <- 95 + 1.8 * bmi + rnorm(35, 0, 8)
\end{rcode}

\textbf{Tasks:}
\begin{enumerate}[label=\alph*)]
\item Calculate correlation and test significance
\item Fit linear regression
\item Report the regression equation
\item What is the predicted SBP for BMI = 30?
\item Create a plot with regression line and 95\% confidence bands
\end{enumerate}

\vspace{2cm}

\subsection{Problem 3: Comparing Correlations}

Compare the correlation between height and weight for males vs. females.

\begin{rcode}
# Data
set.seed(333)
height_male <- rnorm(25, 178, 8)
weight_male <- 0.9 * height_male - 75 + rnorm(25, 0, 6)

height_female <- rnorm(25, 165, 7)
weight_female <- 0.85 * height_female - 65 + rnorm(25, 0, 5)
\end{rcode}

\textbf{Tasks:}
\begin{enumerate}[label=\alph*)]
\item Calculate correlation for males
\item Calculate correlation for females
\item Compare the two correlations. Are they similar?
\item Create side-by-side scatter plots
\item Fit separate regression models for each group
\end{enumerate}

\vspace{2cm}

\section{Saving Your Work}

\textbf{IMPORTANT:} Save your R script!

\begin{rcode}
# At the top of your R script, add:
# Name: Your Name
# Date: Today's Date
# Lab: Week 9 - Correlation and Simple Linear Regression

# Save your workspace
save.image("week9_lab.RData")

# To load it later:
# load("week9_lab.RData")
\end{rcode}

\section{\textcolor{aiblue}{AI-Enhanced Learning Tips}}

\begin{tcolorbox}[colback=aiblue!10,colframe=aiblue,title=Using AI Tools for Regression Analysis]

\textbf{Good prompts for AI assistance:}
\begin{itemize}
\item ``Explain the difference between correlation and regression with examples''
\item ``How do I interpret a regression slope of 1.25 mmHg per year?''
\item ``Create R code to check regression assumptions with diagnostic plots''
\item ``What does it mean if residuals show a curved pattern?''
\item ``Generate scatter plot with regression line and prediction bands''
\end{itemize}

\textbf{AI can help with:}
\begin{itemize}
\item Interpreting regression output
\item Creating publication-quality plots
\item Checking assumptions
\item Troubleshooting R code errors
\item Understanding confidence vs. prediction intervals
\item Explaining statistical concepts
\end{itemize}

\textbf{Remember:}
\begin{itemize}
\item Always verify AI-generated results
\item Check that interpretations make biological sense
\item Examine plots visually, don't just trust numbers
\item AI is a learning tool, you are the analyst!
\end{itemize}
\end{tcolorbox}

\section{Key Takeaways}

\begin{enumerate}
\item \textbf{Correlation:} Measures strength and direction of linear relationship ($-1$ to $+1$)
\item \textbf{Regression:} Models relationship and allows prediction
\item \textbf{Always visualize:} Scatter plots reveal patterns that statistics might miss
\item \textbf{Check assumptions:} Use residual plots (linearity, normality, constant variance)
\item \textbf{Interpret slope in context:} What does 1-unit change in X mean for Y?
\item \textbf{$R^2$:} Proportion of variance explained (0 to 1)
\item \textbf{Prediction intervals:} Wider than confidence intervals (account for individual variation)
\item \textbf{Correlation $\neq$ causation:} Need experiments for causal claims
\end{enumerate}

\section{R Functions Reference}

\begin{itemize}
\item \textbf{Correlation:}
  \begin{itemize}
  \item \texttt{cor(x, y)} -- Correlation coefficient
  \item \texttt{cor.test(x, y)} -- Test significance, get CI
  \item \texttt{cor(x, y, method="spearman")} -- Spearman correlation
  \end{itemize}

\item \textbf{Regression:}
  \begin{itemize}
  \item \texttt{lm(y \textasciitilde\ x)} -- Fit linear model
  \item \texttt{summary(model)} -- Detailed output
  \item \texttt{coef(model)} -- Extract coefficients
  \item \texttt{confint(model)} -- Confidence intervals for coefficients
  \item \texttt{anova(model)} -- ANOVA table
  \end{itemize}

\item \textbf{Prediction:}
  \begin{itemize}
  \item \texttt{predict(model, newdata)} -- Point predictions
  \item \texttt{predict(model, newdata, interval="confidence")} -- CI for mean
  \item \texttt{predict(model, newdata, interval="prediction")} -- PI for individual
  \end{itemize}

\item \textbf{Diagnostics:}
  \begin{itemize}
  \item \texttt{plot(model)} -- Diagnostic plots (residuals, Q-Q, etc.)
  \item \texttt{residuals(model)} -- Extract residuals
  \item \texttt{fitted(model)} -- Extract fitted values
  \item \texttt{shapiro.test(residuals(model))} -- Test normality
  \end{itemize}

\item \textbf{Plotting:}
  \begin{itemize}
  \item \texttt{plot(x, y)} -- Scatter plot
  \item \texttt{abline(model)} -- Add regression line
  \item \texttt{lines(x, y)} -- Add lines to existing plot
  \end{itemize}
\end{itemize}

\section{Additional Resources}

\begin{itemize}
\item R Documentation: \texttt{?lm}, \texttt{?cor.test}, \texttt{?predict}
\item Textbook: Pagano \& Gauvreau, Chapters 17 and 18
\item Quick-R: Linear Regression (\url{https://www.statmethods.net/stats/regression.html})
\item R for Data Science: \url{https://r4ds.had.co.nz/}
\end{itemize}

\end{document}
